% ============================================================
%  Project Gantt Chart — Sprint Timeline
%  Psych Platform — Planning Section
% ============================================================

\section{Project Gantt Chart (Sprint Timeline)}
\label{sec:gantt-sprints}

This section summarises the project plan as a sprint-based timeline. The plan is
organised into five sprints, each with a fixed duration and a bounded scope that
maps to the product backlog (PB) identifiers. The timeline also embeds the
\emph{Definition of Done (DoD)} as the shared completion standard for all sprint
items.

\subsection{Sprint Overview}
\begin{itemize}
  \item \textbf{Sprint 1 (2 weeks)} — Core infrastructure: authentication, role
  system, and base MERN architecture. \textbf{Backlog:} PB-01, PB-02.
  \item \textbf{Sprint 2 (2 weeks)} — Psychologist onboarding: profile creation,
  credential upload, AI verification, and admin approval flow.
  \textbf{Backlog:} PB-03, PB-04, PB-05.
  \item \textbf{Sprint 3 (2 weeks)} — Patient engagement: search and discovery,
  geospatial search, calendar management, and full booking lifecycle.
  \textbf{Backlog:} PB-06 through PB-09.
  \item \textbf{Sprint 4 (3 weeks)} — Core consultation engine: payment and
  session code system, real-time messaging (Socket.IO), voice messages, session
  lifecycle, PDF report generation, and AI chatbot integration.
  \textbf{Backlog:} PB-10 through PB-15.
  \item \textbf{Sprint 5 (2 weeks)} — Platform enrichment: psychologist
  dashboard, document AI querying, ratings, admin panel, i18n (EN/FR/AR + RTL),
  floating assistant, and notifications. \textbf{Backlog:} PB-16 through PB-22.
\end{itemize}

\subsection{Definition of Done (DoD)}
A backlog item is considered \emph{Done} when: (i) all acceptance criteria are
met and verified; (ii) the feature is integrated into the main branch without
regressions; (iii) the relevant API endpoint returns the expected HTTP status
codes under test; and (iv) the UI component renders correctly on both desktop
and mobile viewports.

\subsection{Gantt Chart (Sprints)}
Figure~\ref{fig:gantt-sprints} provides a high-level sprint Gantt chart.
The chart is intentionally sprint-granular (not task-granular) to keep planning
readable and to match the delivery cadence.

\begin{figure}[H]
  \centering
  % Requires: \usepackage{pgfgantt}
  \begin{ganttchart}[
    hgrid,
    vgrid,
    time slot format=isodate-yearmonthday,
    x unit=0.28cm,
    y unit chart=0.6cm,
    bar height=0.5,
    group height=0.3,
    bar/.style={fill=black!20, draw=black!60},
    group/.style={fill=black!10, draw=black!60},
    milestone/.style={fill=black!60, draw=black!60},
    title/.style={draw=none, fill=none},
    title label font=\small\bfseries,
    bar label font=\small,
  ]{2026-01-05}{2026-03-16}
    \gantttitle{Sprint Timeline (Example 11-week Plan)}{71} \\
    \gantttitlecalendar{year, month=name} \\

    \ganttgroup{Sprint 1 (PB-01..PB-02)}{2026-01-05}{2026-01-18} \\
    \ganttbar{Core infrastructure (Auth/Roles/MERN)}{2026-01-05}{2026-01-18} \\

    \ganttgroup{Sprint 2 (PB-03..PB-05)}{2026-01-19}{2026-02-01} \\
    \ganttbar{Psychologist onboarding + admin approval}{2026-01-19}{2026-02-01} \\

    \ganttgroup{Sprint 3 (PB-06..PB-09)}{2026-02-02}{2026-02-15} \\
    \ganttbar{Discovery + geo search + booking lifecycle}{2026-02-02}{2026-02-15} \\

    \ganttgroup{Sprint 4 (PB-10..PB-15)}{2026-02-16}{2026-03-08} \\
    \ganttbar{Consultation engine + realtime + chatbot}{2026-02-16}{2026-03-08} \\

    \ganttgroup{Sprint 5 (PB-16..PB-22)}{2026-03-09}{2026-03-22} \\
    \ganttbar{Enrichment (dashboard/i18n/assistant/notifs)}{2026-03-09}{2026-03-22} \\

    \ganttmilestone{Release candidate}{2026-03-23} \\
  \end{ganttchart}
  \caption{Project Gantt chart: sprint-level delivery timeline (illustrative dates).}
  \label{fig:gantt-sprints}
\end{figure}

The dates shown are illustrative. If a different project start date is required,
the chart can be shifted without changing the sprint durations (2, 2, 2, 3, 2
weeks).

